\chapter{Results and Discussion}

This chapter presents a quantitative analysis of both the experimentally observed droplet trajectories and the trajectories generated by the coupled PDE--SDE model. Rather than treating the experimental and numerical results separately, the comparison is performed through a common set of observables that characterize the dynamical behavior of the system. These observables include the trajectory morphology, mean squared displacement, velocity autocorrelation function, orientational dynamics, and power spectral density. Such an approach enables a direct comparison between the experimental system and the theoretical model.

\section{Trajectory Analysis}

\subsection{Experimental Trajectories}

Figures XX--XX show the trajectories extracted from the experimental image sequences for droplets of different initial volumes. The trajectories exhibit persistent motion, curved paths, and a tendency to avoid previously explored regions. Such behavior differs markedly from classical Brownian motion, where successive displacements are uncorrelated.

The trajectories reveal that the droplets undergo extended periods of directed motion interspersed with reorientation events. This persistence suggests the existence of an underlying memory mechanism influencing the dynamics.

\subsection{PDE--SDE Model Trajectories}

Figures XX--XX show representative trajectories obtained from the coupled PDE--SDE model. Similar to the experimental observations, the simulated droplets exhibit persistent motion and self-avoidance. The trajectories are shaped by the interaction between the stochastic motion of the droplets and the evolving concentration field.

\subsection{Experimental--Model Comparison}

The qualitative agreement between the experimental and simulated trajectories suggests that the model captures the dominant physical mechanisms governing the droplet motion. In particular, both systems display persistence, self-repulsion, and non-random exploration of space.

\section{Mean Squared Displacement Analysis}

The mean squared displacement (MSD) was used to characterize the transport properties of the droplets.

\begin{equation}
\mathrm{MSD}(\tau) = \left\langle |\mathbf r(t+\tau)-\mathbf r(t)|^2 \right\rangle.
\end{equation}

\subsection{Experimental MSD}

Figures XX--XX show the experimentally measured MSDs plotted on logarithmic scales. The MSD exhibits power-law behavior,

\begin{equation}
\mathrm{MSD}(\tau)
\propto
\tau^\alpha,
\end{equation}

where $\alpha$ denotes the transport exponent.

Values of $\alpha > 1$ indicate superdiffusive transport and reveal persistent motion over multiple timescales.

\subsection{Model MSD}

The MSD was computed from the simulated trajectories using the same procedure. The resulting scaling exponents were extracted from power-law fits and compared directly with the experimental values.

\subsection{Comparison}

Table~\ref{tab:msdcomparison} summarizes the MSD exponents obtained from both experiment and simulation.

\begin{table}[h]
\centering
\caption{Comparison of MSD scaling exponents.}
\label{tab:msdcomparison}
\begin{tabular}{lcc}
\hline
Case & $\alpha$ & Transport Regime \\
\hline
Experiment & -- & Superdiffusive \\
PDE--SDE Model & -- & Superdiffusive \\
\hline
\end{tabular}
\end{table}

The agreement between the experimentally measured and simulated MSD exponents indicates that the model successfully reproduces the observed transport properties.

\section{Velocity Autocorrelation Analysis}

The velocity autocorrelation function (VACF) was used to quantify memory effects and directional persistence.

\begin{equation}
C_v(\tau) = \left\langle \mathbf v(t)\cdot \mathbf v(t+\tau) \right\rangle.
\end{equation}

\subsection{Experimental VACF}

The experimentally measured VACFs remain positive over extended lag times and decay gradually rather than collapsing rapidly to zero.

This behavior indicates that the droplets retain memory of their previous direction of motion and therefore exhibit persistent transport.

\subsection{Model VACF}

The VACFs obtained from the PDE--SDE simulations display similar long-lived correlations, indicating that the model captures the persistence observed experimentally.

\subsection{Correlation Time}

The velocity correlation time was computed as

\begin{equation}
\tau_c = \int_0^{\infty} \frac{C_v(\tau)}{C_v(0)} \, d\tau.
\end{equation}

Table~\ref{tab:vacfcomparison} summarizes the extracted correlation times.

\begin{table}[h]
\centering
\caption{Comparison of velocity correlation times.}
\label{tab:vacfcomparison}
\begin{tabular}{lcc}
\hline
Case & $\tau_{c,x}$ & $\tau_{c,y}$ \\
\hline
Experiment & -- & -- \\
PDE--SDE Model & -- & -- \\
\hline
\end{tabular}
\end{table}

The persistence observed in the VACF provides direct evidence that the droplet dynamics cannot be adequately described by a memoryless random walk.

\section{Orientational Dynamics}

The orientation of motion was characterized using

\begin{equation}
\theta(t) = \tan^{-1}\left( \frac{v_y(t)}{v_x(t)} \right).
\end{equation}

The orientational correlation functions reveal finite directional persistence in both experiment and simulation. The droplets maintain their direction of motion over significant time intervals before undergoing reorientation events.

The observed orientational persistence is consistent with both the MSD and VACF analyses.

\section{Power Spectral Density Analysis}

The power spectral density (PSD) of the velocity components was computed to investigate the frequency content of the motion.

The spectra exhibit enhanced low-frequency power, indicating the presence of long-timescale dynamical correlations. Similar behavior is observed in both experiment and simulation.

The PSD therefore provides complementary evidence for persistent motion and memory effects.

\section{Summary of Dynamical Quantities}

Table~\ref{tab:summary} summarizes the principal observables extracted from both the experiments and the PDE--SDE model.

\begin{table}[h]
\centering
\caption{Summary of experimentally measured and simulated dynamical quantities.}
\label{tab:summary}
\begin{tabular}{lcc}
\hline
Observable & Experiment & PDE--SDE Model \\
\hline
MSD exponent $\alpha$ & -- & -- \\
Correlation time $\tau_{c,x}$ & -- & -- \\
Correlation time $\tau_{c,y}$ & -- & -- \\
Mean speed $\langle v\rangle$ & -- & -- \\
Orientation persistence & -- & -- \\
Low-frequency PSD enhancement & Yes & Yes \\
\hline
\end{tabular}
\end{table}

\section{Physical Interpretation}

The combined experimental and numerical results support a common physical interpretation. The droplets behave as active self-propelled particles coupled to a slowly evolving chemical environment.

The MSD analysis reveals superdiffusive transport, while the VACF demonstrates long-lived velocity correlations. The orientational correlation functions further indicate directional persistence, and the PSD analysis reveals the existence of long-timescale dynamical modes.

Together, these observations indicate that the droplets retain memory of their previous motion.

A plausible physical mechanism for this memory is the chemical trail generated by the droplets. As the droplet moves, it modifies the surrounding concentration field through dissolution and chemical release. The concentration field evolves diffusively and subsequently influences the future propulsion of the droplet.

This feedback loop can be summarized as follows:

\begin{enumerate}
\item the droplet moves and generates a chemical trail,
\item the trail modifies the local concentration field,
\item concentration gradients alter the propulsion force,
\item future motion depends on previous trajectories.
\end{enumerate}

Consequently, the dynamics are effectively non-Markovian. The future evolution of the droplet cannot be determined solely from its instantaneous state because information about previous motion is stored in the surrounding chemical field.

The ability of the PDE--SDE model to reproduce the experimentally observed persistence, superdiffusion, and self-avoidance strongly supports this interpretation and suggests that chemical-memory-driven feedback is the dominant mechanism governing the droplet dynamics.
